\section{USAMO 2007}
        \begin{enumerate}
        	\item (\textit{USAMO 2007}) Let $n$ be a positive integer. Define a sequence by setting $a_1=n$ and, for each $k>1$, letting $a_k$ be the unique integer in the range $0\le a_k\le k-1$ for which $a_1+a_2+\cdots+a_k$ is divisible by $k$. For instance, when $n=9$ the obtained sequence is $9, 1, 2, 0, 3, 3, 3, \ldots$.  Prove that for any $n$ the sequence $a_1, a_2, a_3, \ldots$ eventually becomes constant.


 



\textbf{Solution 1}

Let $S_k = a_1 + a_2 + \cdots + a_k$ and $b_k = \frac{S_k}{k}$. Thus, because $S_{k+1} = S_k + a_{k+1}$,
 \[b_{k+1} = \frac{b_k \cdot k + a_{k+1}}{k+1} = \left(\frac{k}{k+1}\right) \cdot b_k + \frac{a_{k+1}}{k+1}\]$\frac{k}{k+1} < 1$, and by definition, $\frac{a_{k+1}}{k+1} < 1$. Thus, $b_{k+1} < b_k + 1$. Also, both $b_k$ and $b_{k+1}$ are integers, so $b_{k+1} \le b_k$. As the $b_k$'s form a \href{https://artofproblemsolving.com/wiki/index.php?title=Non-increasing&action=edit&redlink=1}{non-increasing} sequence of positive integers, they must eventually become constant. 

 Therefore, $b_k = b_{k+1}$ for some sufficiently large value of $k$. Then $a_{k+1} = S_{k+1} - S_k = b_k(k + 1) - b_k(k) = b_k$, so eventually the sequence $a_k$ becomes constant.

 



\textbf{Solution 2}

Let $a_1 = n$. Since $a_k\leq k - 1$, we have that
 \[a_1 + a_2 + \cdots + a_n\leq n + 1 + 2 + \cdots + n - 1 = \frac{n(n + 1)}{2}.\]Since $a_1 + a_2 + \cdots + a_n = nk$ for some integer $k$, we can keep adding $k$ to satisfy the conditions, provided that $k\leq n$. This is true since $k\leq\frac{n + 1}{2}\leq n$, so the sequence must eventually become constant.

 



\textbf{Solution 3}

Define $S_k = a_1 + a_2 + \cdots + a_k$, and $b_k = \frac{S_k}{k}$. By the problem hypothesis, $b_k$ is an integer valued sequence.

 \textbf{Lemma:} There exists a $k$ such that $b_k < k$.

 \textit{Proof:} Choose any $k$ such that $k^2 + 3k - 2 > 2n$. Then
 \begin{align*} \frac{k^2 + 3k - 2}{2} &> n \\ k^2 &> \frac{k^2 - 3k + 2}{2} + n \\ k^2 &> (k-2) + (k-1) + \cdots + 1 + n \\ k^2 &> a_{k-1} + a_{k-2} + \cdots + a_2 + a_1 \\ k^2 &> S_k \\ k &> \frac{S_k}{k} \\ k &> b_k, \end{align*}
as desired.

 \textbf{End Lemma}

 Let $k$ be the smallest $k$ such that $b_k < k$. Then $b_k = m < k$, and $S_k = km$. To make $b_{k+1}$ an integer, $S_{k+1} = S_k + a_{k+1}$ must be divisible by $k+1$. Thus, because $km + m$ is divisible by $k+1$, $a_{k+1} \equiv m \pmod{k+1}$, and, because $0 \le a_{k+1} < k$, $a_{k+1} = m$. Then $b_{k+1} = \frac{(k+1)m}{k+1} = m$ as well. Repeating the same process using $k+1$ instead of $k$ gives $a_{k+2} = m$, and an easy induction can prove that for all $N > k+1$, $a_N = m$. Thus, $a_k$ becomes a constant function for arbitrarily large values of $k$.

 



\textbf{Solution 4}

For $k\geq 1$, let
 \[s_k = a_1 + a_2 + \cdots + a_k.\]We claim that for some $m$ we have $s_m = m(m - 1)$. To this end, consider the sequence which computes the differences between $s_k$ and $k(k - 1)$, i.e., whose $k$-th term is $s_k - k(k - 1)$. Note that the first term of this sequence is positive (it is equal to $n$) and that its terms are strictly decreasing since
 \[(s_k - k(k - 1)) - (s_{k+1} - (k + 1)k) = 2k - a_{k+1}\geq 2k - k = k\geq 1.\]Further, a negative term cannot immediately follow a positive term. Suppose otherwise, namely that $s_k > k(k - 1)$ and $s_{k+1} < (k + 1)k$. Since $s_k$ and $s_{k+1}$ are divisible by $k$ and $k + 1$, respectively, we can tighten the above inequalities to $s_k\geq k^2$ and $s_{k+1}\leq (k + 1)(k - 1) = k^2 - 1$. But this would imply that $s_k > s_{k+1}$, a contradiction. We conclude that the sequence of differences must eventually include a term equal to zero.

 Let $m$ be a positive integer such that $s_m = m(m - 1)$. We claim that
 \[m - 1 = a_{m+1} = a_{m+2} = a_{m+3} = a_{m+4} = \cdots.\]This follows from the fact that the sequence $a_1, a_2, a_3, \ldots$ is uniquely determined and choosing $a_{m+i} = m - 1$, for $i\geq 1$, satisfies the range condition
 \[0\leq a_{m+i} = m - 1\leq m + i - 1,\]and yields
 \[s_{m+i} = s_m + i(m - 1) = m(m - 1) + i(m - 1) = (m + i)(m - 1).\]

 


 



\textbf{Solution 5}

Let $S_k =\sum_{i=1}^k a_i$, where $S_k \equiv 0 \pmod{k}$ and $0 \leq a_k < k$. Define $m_k = S_k / k$, with $m_1 = n$.

 First, for $k > 1$, since $S_{k-1} = (k-1) \cdot m_{k-1}$ and $S_k = S_{k-1} + a_k$, we need $S_k \equiv 0 \pmod{k}$. Set $a_k \equiv m_{k-1} \pmod{k}$. Then,

 $S_k \bmod k \equiv \bigl((k-1) \cdot m_{k-1} + (m_{k-1} \bmod k)\bigr) \pmod{k}$.Since $(k-1) \cdot m_{k-1} \equiv -m_{k-1} \pmod{k}$, and $-m_{k-1} + (m_{k-1} \bmod k) \equiv 0 \pmod{k}$, this holds.

 Next, compute $m_k = S_k / k = \bigl((k-1) \cdot m_{k-1} + a_k\bigr) / k$. If $m_{k-1} \geq k$, let $m_{k-1} = qk + r$ with $0 \leq r < k$, so $a_k = r$, and$m_k = \frac{(k-1) \cdot m_{k-1} + r}{k} < m_{k-1}$,so $m_k < m_{k-1}$. Moreover, so long as $m_{k-1} \geq k$ then $m_{i}$ is strictly decreasing while $k$ is strictly increasing, until $m_{i} \leq k$, in which case the sequence terminates as $m_{k-1} = m_{k}$

 



\textbf{Solution 6(like solution 2)}

First, we may make an observation and say that for $n \equiv p (mod k)$, $\sum_{m=2}^{k} a_m \equiv k-p (mod k)$ must occur for the whole sum to be divisible by $k$. Thus, the following is apparent:
 \[\sum_{m=2}^{k} a_m =(q+1)k - p , k < n\]Then, we may make another observation that when $n=k$, the sum also has to be divisible by n. We may then explore when n=k:
 \[\sum_{m=2}^{n} a_m \equiv 0 (mod n), \sum_{m=2}^{n} a_m = rn\]and
 \[a_n = \sum_{m=2}^{n} a_m - \sum_{m=2}^{n-1} a_m\]Then,
 \[a_n = rn - \sum_{m=2}^{n-1} a_m \le n-1\]Also,
 \[a_{n+s} = r, r \le n\] for $s \ge 1$. This is because:
 \[\sum_{m=2}^{n} a_m + r = rn + r = r(n+1)\]This must be true since $r(n+1)$ will be divisible by $n+1$ and $k=n+1$, we may then generalize this to all $r(n+s), s \in \mathbb{Z}, s \ge 1$
 \[rn + sr= r(n+s), \frac{r(n+s)}{r} = n + s = \frac{\sum_{m=2}^{n+s} a_m}{r}\]Thus, we may say that the sequence $a_1,a_2,...a_k$ must converge to some integer value $r \le n$ when $k \ge n + 1$.

 



\textbf{Solution 7(A simpler version of Solution 1)}

$a_1=1 \cdot x_1$, $a_1+a_2=2 \cdot x_2$, $a_1+a_2+a_3=3 \cdot x_3, ...$

 Claim.

 $x_1 \geq x_2 \geq x_3 \geq ...$

 Proof.

 $a_2=2 \cdot x_2 - 1 \cdot x_1$, $a_3=3 \cdot x_3 - 2 \cdot x_2$, $a_4=4 \cdot x_4 - 3 \cdot x_3, ...$

 F.T.S.O.C suppose $x_{n-1} < x_n$ for some $n>1$.

 $a_n=n \cdot x_n - (n-1) \cdot x_{n-1}$.

 Note that $n \cdot x_n > (n-1) \cdot (x_{n-1}+1)$ as $x_n \geq x_{n-1}+1$.

 This directly implies $n \cdot x_n - (n-1) \cdot x_{n-1} = a_n > n-1$, a contradiction to the bound given on $a_n$ by the problem.

 Hence, it is impossibe to have $x_{n-1} < x_n$ for any $n>1$, implying  $x_{n-1} \geq x_n$ for each $n>1$, that is, $x_1 \geq x_2 \geq x_3 \geq ...$, proving our claim.

 
As the sequence $x_i$ only consists of positive integers, it can't be ever decreasing, so that it must become eventually constant, 

 that is, $x_k=x_{k+1}=...$ from some $k \geq 1$.

 Then, 

 $a_1+...+a_k=k \cdot x_k$, $a_1+...+a_k+a_{k+1}=(k+1) \cdot x_{k+1}=(k+1) \cdot x_k, ...$

 giving us $x_k=a_{k+1}=a_{k+2}=...$ (by substracting the 1st equation from the 2nd, the 2nd from the 3rd, and so on), completing the proof to the problem.

 <i>Alternate solutions are always welcome. If you have a different, elegant solution to this problem, please \href{https://artofproblemsolving.com/wiki/index.php?title=2007\_USAMO\_Problems/Problem\_1&action=edit}{add it} to this page.</i>

 


	\item (\textit{USAMO 2007}) A square grid on the Euclidean plane consists of all points $(m,n)$, where $m$ and $n$ are integers. Is it possible to cover all grid points by an infinite family of discs with non-overlapping interiors if each disc in the family has radius at least 5?


 



\textbf{Solution 1}

\textbf{Lemma.} Among 3 \href{https://artofproblemsolving.com/wiki/index.php?title=Tangent}{tangent} circles with radius greater than or equal to 5, one can always fit a circle with radius greater than $\frac{1}{\sqrt{2}}$ between those 3 circles. 

 \textit{Proof.} \href{https://artofproblemsolving.com/wiki/index.php?title=Descartes\%27\_Circle\_Theorem&action=edit&redlink=1}{Descartes' Circle Theorem} states that if $a$ is the curvature of a circle ($a=\frac 1{r}$, positive for \href{https://artofproblemsolving.com/wiki/index.php?title=Externally\_tangent}{externally tangent}, negative for \href{https://artofproblemsolving.com/wiki/index.php?title=Internally\_tangent}{internally tangent}), then we have that 
 \[(a+b+c+d)^2=2(a^2+b^2+c^2+d^2)\]Solving for $a$, we get 
 \[a=b+c+d+2 \sqrt{bc+cd+db}\]Take the positive root, as the negative root corresponds to internally tangent circle.

 Now clearly, we have $b+c+d \le \frac 35$, and $bc+cd+db\le \frac 3{25}$.Summing/\href{https://artofproblemsolving.com/wiki/index.php?title=Square\_root}{square root}/multiplying appropriately shows that $a \le \frac{3 + 2 \sqrt{3}}5$. Incidently, $\frac{3 + 2\sqrt{3}}5 < \sqrt{2}$, so $a< \sqrt{2}$, $r > \frac 1{\sqrt{2}}$, as desired. $\blacksquare$

 For sake of \href{https://artofproblemsolving.com/wiki/index.php?title=Contradiction}{contradiction}, assume that we have a satisfactory placement of circles. Consider 3 circles, $p,\ q,\ r$ where there are no circles in between. By \href{https://artofproblemsolving.com/wiki/index.php?title=Appolonius\%27\_problem&action=edit&redlink=1}{Appolonius' problem}, there exists a circle $t$ tangent to $p,\ q,\ r$ externally that is between those 3 circles. Clearly, if we move $p,\ q,\ r$ together, $t$ must decrease in radius. Hence it is sufficient to consider 3 tangent circles. By lemma 1, there is always a circle of radius greater than $\frac{1}{\sqrt{2}}$ that lies between $p,\ q,\ r$. However, any circle with $r>\frac 1{\sqrt{2}}$ must contain a \href{https://artofproblemsolving.com/wiki/index.php?title=Lattice\_point}{lattice point}. (Consider placing an unit square parallel to the gridlines in the circle.) That is a contradiction. Hence no such tiling exists.

 



\textbf{Solution 2}

It is not possible. The proof is by contradiction. Suppose that such a covering family $\mathcal{F}$ exists. Let $D(P,\rho)$ denote the disc with center $P$ and radius $\rho$. Start with an arbitrary disc $D(O,r)$ that does not overlap any member of $\mathcal{F}$. Then $D(O,r)$ covers no grid point. Take the disc $D(O,r)$ to be maximal in the sense that any further enlargement would cause it to violate the non-overlap condition. Then $D(O,r)$ is tangent to at least three discs in $\mathcal{F}$. Observe that there must be two of the three tangent discs, say $D(A,a)$ and $D(B,b)$ such that $\angle AOB\leq 120^\circ$. By the Law of Cosines applied to triangle $ABO$,
 \[(a + b)^2\leq (a + r)^2 + (b + r)^2 + (a + r)(b + r),\]which yields
 \[ab\leq 3(a + b)r + 3r^2,\]and thus
 \[12r^2\geq (a - 3r)(b - 3r).\]Note that $r < 1/\sqrt{2}$ because $D(O,r)$ covers no grid point, and $(a - 3r)(b - 3r)\geq (5 - 3r)^2$ because each disc in $\mathcal{F}$ has radius at least 5. Hence $2\sqrt{3}r\geq 5 - 3r$, which gives $5\leq (3 + 2\sqrt{3})r < (3 + 2\sqrt{3})/\sqrt{2}$ and thus $5\sqrt{2} < 3 + 2\sqrt{3}$. Squaring both sides of this inequality yields $50 < 21 + 12\sqrt{3} < 21 + 12\cdot 2 = 45$. This contradiction completes the proof.

 \textbf{Remark:} The above argument shows that no covering family exists where each disc has radius greater than $(3 + 2\sqrt{3})/\sqrt{2}\approx 4.571$. In the other direction, there exists a covering family in which each disc has radius $\sqrt{13}/2\approx 1.802$. Take discs with this radius centered at points of the form $\left(2m + 4n + \frac{1}{2}, 3m + \frac{1}{2}\right)$, where $m$ and $n$ are integers. Then any grid point is with $\sqrt{13}/2$ of one of the centers and the distance between any two centers is at least $\sqrt{13}$. The extremal radius of a covering family is unknown.

 <i>Alternate solutions are always welcome. If you have a different, elegant solution to this problem, please \href{https://artofproblemsolving.com/wiki/index.php?title=2007\_USAMO\_Problems/Problem\_2&action=edit}{add it} to this page.</i>

 


	\item (\textit{USAMO 2007}) Let $S$ be a set containing $n^2+n-1$ elements, for some positive integer $n$. Suppose that the $n$-element subsets of $S$ are partitioned into two classes. Prove that there are at least $n$ pairwise disjoint sets in the same class.


 



\textbf{Solution 1}

Claim: If we have instead $k(n+1)-1$ elements, then we must have k disjoint subsets in a class.We proceed by Induction

 Base Case: $k=1$ is trivial.

 Case 1: If there exists an $n+1$ element subset $X$ s.t. $X$ has 2 n-element subsets in different classes.Consider the $(k-1)(n+1)-1$ elements that are not in $X$, from the Induction hypothesis, we have $k-1$ disjoint subsets in one of the classes. Those $k-1$ subsets paired with a subset of $X$ in the same class forms $k$ disjoint subsets of size n.

 Case 2: All $n+1$ element subsets have all of their n-element subsets in one class.Look at any 2 n-element subsets $a_1$ and $a_p$. Then we can create a path $a_1,a_2,...,a_p$ s.t. any 2 consecutive terms belong to the same n+1-element subset.(i.e. we change 1 element each time.) Thus, $a_1$ and $a_p$ are in the same class and we see that all subsets are in one class. Obviously we can find $k$ disjoint subsets of size n.

 
The question is asking for the case when we substitute $k=n$ which we have proved.

 



\textbf{Solution 2}

Call an $n+1$-element subset of $S$ separable if it has a subset in each class of the partition. We \href{https://artofproblemsolving.com/wiki/index.php?title=Recursion}{recursively} build a set $Q$ of disjoint separable subsets of $S$: begin with $Q$ empty and at each step if there is a separable subset which is disjoint from all sets in $Q$ add that set to $Q$. The process terminates when every separable subset intersects a set in $Q$. Let $T$ be the set of elements in $S$ which are not in any set in $Q$. We claim that one class contains every $n$-element subset of $T$.

 Suppose that $a_1, a_2, \ldots a_k$ are elements of $T$. Denote by $A_i$ the set $\left\{a_i, a_{i+1}, \ldots a_{i+n-1}\right\}$. Note that for each $i$, $A_i \cup A_{i+1}$ is not separable, so that $A_i$ and $A_{i+1}$ are in the same class. But then $A_i$ is in the same class for each $1 \leq i \leq k - n + 1$ — in particular, $A_1$ and $A_{k-n+1}$ are in the same class. But for any two sets we may construct such a sequence with $A_1$ equal to one and $A_{k-n+1}$ equal to the other.  

 We are now ready to construct our $n$ disjoint sets. Suppose that $|Q| = q$. Then $|T| = (n+1)(n-q) - 1 \geq n(n-q)$, so we may select $n - q$ disjoint $n$-element subsets of $T$. Then for each of the $q$ sets in $Q$, we may select a subset which is in the same class as all the subsets of $T$, for a total of $n$ disjoint sets.

 



\textbf{Solution 3}

In order to apply induction, we generalize the result to be proved so that it reads as follows:

 \textbf{Proposition.} If the $n$-element subsets of a set $S$ with $(n + 1)m - 1$ elements are partitioned into two classes, then there are at least $m$ pairwise disjoint sets in the same class.

 \textit{Proof.} Fix $n$ and proceed by induction on $m$. The case of $m = 1$ is trivial. Assume $m > 1$ and that the proposition is true for $m - 1$. Let $\mathcal{P}$ be the partition of the $n$-element subsets into two classes. If all the $n$-element subsets belong to the same class, the result is obvious. Otherwise select two $n$-element subsets $A$ and $B$ from different classes so that their intersection has maximal size. It is easy to see that $|A\cap B| = n - 1$. (If $|A\cap B| = k < n - 1$, then build $C$ from $B$ by replacing some element not in $A\cap B$ with an element of $A$ not already in $B$. Then $|A\cap C| = k + 1$ and $|B\cap C| = n - 1$ and either $A$ and $C$ or $B$ and $C$ are in different classes.) Removing $A\cup B$ from $S$, there are $(n + 1)(m - 1) - 1$ elements left. On this set the partition induced by $\mathcal{P}$ has, by the inductive hypothesis, $m - 1$ pairwise disjoint sets in the same class. Adding either $A$ or $B$ as appropriate gives $m$ pairwise disjoint sets in the same class. $\blacksquare$

 \textbf{Remark:} The value $n^2 + n - 1$ is sharp. A set $S$ with $n^2 + n - 2$ elements can be split into a set $A$ with $n^2 - 1$ elements and a set $B$ of $n - 1$ elements. Let one class consist of all $n$-element subsets of $A$ and the other consist of all $n$-element subsets that intersect $B$. Then neither class contains $n$ pairwise disjoint sets.

 


 <i>Alternate solutions are always welcome. If you have a different, elegant solution to this problem, please \href{https://artofproblemsolving.com/wiki/index.php?title=2007\_USAMO\_Problems/Problem\_3&action=edit}{add it} to this page.</i>

 


	\item (\textit{USAMO 2007}) An \textit{animal} with $n$ \textit{cells} is a connected figure consisting of $n$ equal-sized cells.${}^1$ The figure below shows an 8-cell animal.


 
                    \begin{figure}[H]
                        \centering
                        \includegraphics[width=0.5\textwidth]{Problem/USAMO_I_2007_4_Problem_Figure_1.PNG}
                    \end{figure}
                    A \textit{dinosaur} is an animal with at least 2007 cells. It is said to be \textit{primitive} if its cells cannot be partitioned into two or more dinosaurs. Find with proof the maximum number of cells in a primitive dinosaur.


 <small>${}^1$Animals are also called \textit{polyominoes}. They can be defined inductively. Two cells are \textit{adjacent} if they share a complete edge. A single cell is an animal, and given an animal with $n$ cells, one with $n+1$ cells is obtained by adjoining a new cell by making it adjacent to one or more existing cells.</small>


 



\textbf{Solution 1}

Let a $n$-dino denote an animal with $n$ or more cells.

 We show by \href{https://artofproblemsolving.com/wiki/index.php?title=Induction}{induction} that an $n$-dino with $4n-2$ or more animal cells is not primitive. (Note: if it had more, we could just take off enough until it had $4n-2$, which would have a partition, and then add the cells back on.)

 Base Case: If $n=1$, we have two cells, which are clearly not primitive.

 Inductive Step: Assume any $4n-2$ cell animal can be partitioned into two or more $n$-dinos.

 For a given $(4n+2)$-dino, take off any four cells (call them $w,\ x,\ y,\ z$) to get an animal with $4n-2$ cells.

 This can be partitioned into two or more $n$-dinos, let's call them $A$ and $B$. This means that $A$ and $B$ are connected.

 If both $A$ and $B$ are $(n+1)$-dinos or if $w,\ x,\ y,\ z$ don't all attach to one of them, then we're done.

 So assume $A$ has $n$ cells and thus $B$ has at least $3n-2$ cells, and that $w,\ x,\ y,\ z$ are added to $B$. So $B$ has $3n+2$ cells total.

 Let $C$ denote the cell of $B$ attached to $A$. There are $3n+1$ cells on $B$ besides $C$. Thus, of the three (or less) sides of $C$ not attached to $A$, one of them must have $n+1$ cells by the \href{https://artofproblemsolving.com/wiki/index.php?title=Pigeonhole\_principle}{pigeonhole principle}. It then follows that we can add $A$, $C$, and the other two sides together to get an $(n+1)$ dino, and the side of $C$ that has $n+1$ cells is also an $n+1$-dino, so we can partition the animal with $4n+2$ cells into two $(n+1)$-dinos and we're done.

 Thus, our answer is $4(2007) - 3 = 8025$ cells.

 <table><tbody><tr><td width="50\%">
                    \begin{figure}[H]
                        \centering
                        \includegraphics[width=0.5\textwidth]{Solution/USAMO_I_2007_4_Solution_Figure_1.PNG}
                    \end{figure}
                    </td><td width="50\%">
                    \begin{figure}[H]
                        \centering
                        \includegraphics[width=0.5\textwidth]{Solution/USAMO_I_2007_4_Solution_Figure_2.PNG}
                    \end{figure}
                    </td></tr><tr><td width="50\%">Example of a solution</td><td width="50\%">Attempting to partition solution into $\ge 2$ dinosaurs</td></tr></tbody></table>



\textbf{Solution 2}

For simplicity, let $k=2007$ and let $n$ be the number of \href{https://artofproblemsolving.com/wiki/index.php?title=Square}{squares}. Let the centers of the squares be \href{https://artofproblemsolving.com/wiki/index.php?title=Vertices}{vertices}, and connect any centers of adjacent squares with edges. Suppose we have some \href{https://artofproblemsolving.com/wiki/index.php?title=Loop&action=edit&redlink=1}{loops}. Just remove an edge in the loop. We are still connected since you can go around the other way in the loop. Now we have no loops. Each vertex can have at most 4 edges coming out of it. For each point, assign it the \href{https://artofproblemsolving.com/wiki/index.php?title=Quadruple&action=edit&redlink=1}{quadruple}: $(a,b,c,d)$ where $a$, $b$, $c$, $d$ are the numbers of vertices on each branch, \href{https://artofproblemsolving.com/wiki/index.php?title=WLOG}{WLOG} $a\ge b\ge c\ge d$. Note $a+b+c+d=n-1$. 

 Claim: If $n=4k-2$, then we must be able to divide the animal into two dinosaurs. Chose a vertex, $v$, for which $a$ is minimal (i.e. out of all maximal elements in a quadruple, choose the one with the least maximal element).  We have that $4a \ge a+b+c+d=4k-3$, so $a\ge k$. Hence we can just cut off that branch, that forms a dinosaur. 

 But suppose the remaining vertices do not make a dinosaur. Then we have $b+c+d+1\le k-1 \iff n-a\le k-1\iff a\ge 3k-1$. Now move to the first point on the branch at $a$. We have a new quadruple $p,\ q,\ r,\ b+c+d+1$) where $p+q+r=a-1\ge 3k-2$. 

 Now consider the maximal element of that quadruple. We already have $b+c+d+1\le k-1$. WLOG $p\ge q\ge r\ge 0$, then $3p\ge p+q+r=a-1\ge 3k-2\implies p\ge k$ so $p>k-1=b+c+d+1$, so $p$ is the maximal element of that quadruple. 

 Also $a-1=p+q+r\ge p+0+0$, so $p<a$. But that is a \href{https://artofproblemsolving.com/wiki/index.php?title=Contradiction}{contradiction} to the minimality of $a$. Therefore, we must have that $b+c+d+1\ge k$, so we have a partition of two dinosaurs.

 Maximum: $n=4k-3$.Consider a cross with each branch having $k-1$ vertices. Clearly if we take partition $k$ vertices, we remove the center, and we are not connected.

 So $k=2007$: $4\cdot 2007-3=8025$.

 



\textbf{Solution 3 (Generalization)}

Turn the dinosaur into a \href{https://artofproblemsolving.com/wiki/index.php?title=Graph}{graph} (cells are \href{https://artofproblemsolving.com/wiki/index.php?title=Vertex}{vertices}, adjacent cells connected by an edge) and prove this result about graphs. A connected graph with $V$ vertices, where each vertex has degree less than or equal to $D$, can be partitioned into connected components of sizes at least $\frac{V-1}{D}$. So then in this special case, we have $D = 4$, and so $V = 2006 \times 4+1$ (a possible configuration of this size that works consists of a center and 4 lines of cells each of size 2006 connected to the center). We next throw out all the \href{https://artofproblemsolving.com/wiki/index.php?title=Geometry}{geometry} of this situation, so that we have a completely unconstrained graph. If we prove the above-mentioned result, we can put the geometry back in later by taking the connected components that our partition gives us, then filling back all edges that have to be there due to adjacent cells. This won't change any of the problem constraints, so we can legitimately do this. 

 Going, now, to the case of arbitrary graphs, we WOP on the number of edges. If we can remove any edge and still have a connected graph, then we have found a smaller graph that does not obey our theorem, a contradiction due to the minimality imposed by WOP. Therefore, the only case we have to worry about is when the graph is a tree. If it's a tree, we can root the tree and consider the size of subtrees. Pick the root such that the size of the largest subtree is minimized. This minimum must be at least $\frac{V-1}{D}$, otherwise the sum of the size of the subtrees is smaller than the size of the graph, which is a contradiction. Also, it must be at most $\frac{V}{2}$, or else pick the subtree of size greater than $\frac{V}{2}$ and you have decreased the size of the largest subtree if you root from that vertex instead, so you have some subtree with size between $\frac{V-1}{D}$ and $\frac V2$. Cut the edge connecting the root to that subtree, and use that as your partition. 

 It is easy to see that these partitions satisfy the contention of our theorem, so we are done.

 



\textbf{Solution 4}

Let $s$ denote the minimum number of cells in a dinosaur; the number this year is $s = 2007$.

 \textbf{Claim:} The maximum number of cells in a primitive dinosaur is $4(s - 1) + 1$.

 First, a primitive dinosaur can contain up to $4(s - 1) + 1$ cells. To see this, consider a dinosaur in the form of a cross consisting of a central cell and four arms with $s - 1$ cells apiece. No connected figure with at least $s$ cells can be removed without disconnecting the dinosaur.

 The proof that no dinosaur with at least $4(s - 1) + 2$ cells is primitive relies on the following result.

 \textbf{Lemma.} Let $D$ be a dinosaur having at least $4(s - 1) + 2$ cells, and let $R$ (red) and $B$ (black) be two complementary animals in $D$, i.e. $R\cap B = \emptyset$ and $R\cup B = D$. Suppose $|R|\leq s - 1$. Then $R$ can be augmented to produce animals $\tilde{R}\subset R$ and $\tilde{B} = D\backslash\tilde{R}$ such that at least one of the following holds:

 (i) $|\tilde{R}|\geq s$ and $|\tilde{B}|\geq s$,

 (ii) $|\tilde{R}| = |R| + 1$,

 (iii) $|R| < |\tilde{R}|\leq s - 1$.

 \textit{Proof.} If there is a black cell adjacent to $R$ that can be made red without disconnecting $B$, then (ii) holds. Otherwise, there is a black cell $c$ adjacent to $R$ whose removal disconnects $B$. Of the squares adjacent to $c$, at least one is red, and at least one is black, otherwise $B$ would be disconnected. Then there are at most three resulting components $\mathcal{C}_1, \mathcal{C}_2, \mathcal{C}_3$ of $B$ after the removal of $c$. Without loss of generality, $\mathcal{C}_3$ is the largest of the remaining components. (Note that $\mathcal{C}_1$ or $\mathcal{C}_2$ may be empty.) Now $\mathcal{C}_3$ has at least $\lceil (3s - 2)/3\rceil = s$ cells. Let $\tilde{B} = \mathcal{C}_3$. Then $|\tilde{R}| = |R| + |\mathcal{C}_1| + |\mathcal{C}_2| + 1$. If $|\tilde{B}|\leq 3s - 2$, then $|\tilde{R}|\geq s$ and (i) holds. If $|\tilde{B}|\geq 3s - 1$ then either (ii) or (iii) holds, depending on whether $|\tilde{R}|\geq s$ or not. $\blacksquare$

 Starting with $|R| = 1$, repeatedly apply the Lemma. Because in alternatives (ii) and (iii) $|R|$ increases but remains less than $s$, alternative (i) eventually must occur. This shows that no dinosaur with at least $4(s - 1) + 2$ cells is primitive.

 <i>Alternate solutions are always welcome. If you have a different, elegant solution to this problem, please \href{https://artofproblemsolving.com/wiki/index.php?title=2007\_USAMO\_Problems/Problem\_4&action=edit}{add it} to this page.</i>

 


	\item (\textit{USAMO 2007}) Prove that for every nonnegative integer $n$, the number $7^{7^n}+1$ is the product of at least $2n+3$ (not necessarily distinct) primes.


 



\textbf{Solution 1}

The proof is by induction. The base is provided by the $n = 0$ case, where $7^{7^0} + 1 = 7^1 + 1 = 2^3$. To prove the inductive step, it suffices to show that if $x = 7^{2m - 1}$ for some positive integer $m$ then $(x^7 + 1)/(x + 1)$ is composite. As a consequence, $x^7 + 1$ has at least two more prime factors than does $x + 1$. To confirm that $(x^7 + 1)/(x + 1)$ is composite, observe that
 \begin{align*} \frac{x^7 + 1}{x + 1} &= \frac{(x + 1)^7 - ((x + 1)^7 - (x^7 + 1))}{x + 1} \\ &= (x + 1)^6 - \frac{7x(x^5 + 3x^4 + 5x^3 + 5x^2 + 3x + 1)}{x + 1} \\ &= (x + 1)^6 - 7x(x^4 + 2x^3 + 3x^2 + 2x + 1) \\ &= (x + 1)^6 - 7^{2m}(x^2 + x + 1)^2 \\ &= \{(x + 1)^3 - 7^m(x^2 + x + 1)\}\{(x + 1)^3 + 7^m(x^2 + x + 1)\}. \end{align*}
Also each factor exceeds 1. It suffices to check the smaller one; $\sqrt{7x}\leq x$ gives
 \begin{align*} (x + 1)^3 - 7^m(x^2 + x + 1) &= (x + 1)^3 - \sqrt{7x}(x^2 + x + 1) \\ &\geq x^3 + 3x^2 + 3x + 1 - x(x^2 + x + 1) \\ &= 2x^2 + 2x + 1\geq 113 > 1. \end{align*}
Hence $(x^7 + 1)/(x + 1)$ is composite and the proof is complete.

 



\textbf{Solution 2}

We prove by induction on $n$ that $7^{7^n} + 1$ is the product of at least $2n + 3$ (not necessarily distinct) primes.

 For $n = 0$, we have $7^{7^0} + 1 = 7 + 1 = 8 = 2^3$, which has $3 = 2(0) + 3$ prime factors.  Thus, the base case holds.

 Assume that for some $n \ge 0$, the number $7^{7^n} + 1$ is the product of at least $2n + 3$ primes.  We must show that $7^{7^{n+1}} + 1$ has at least $2n + 5$ prime factors.

 Let $a = 7^n$. Since $a$ is odd, we can apply the Aurifeuillian factorization:$7^{7a} + 1 = (7^a + 1)\big(7^{3a} + 3\cdot7^{2a} + 3\cdot7^a + 1\big)\big(7^{\frac{5a}{2} + \frac{1}{2}} + 7^{\frac{3a}{2} + \frac{1}{2}} + 7^{\frac{a}{2} + \frac{1}{2}}\big)$.

 Each factor is an integer greater than 1, so the factorization is nontrivial.

 The first factor is $7^a + 1 = 7^{7^n} + 1$, which by the inductive hypothesis has at least $2n + 3$ prime factors.  It remains to show that the quotient$\frac{7^{7a} + 1}{7^a + 1}$is composite.

 From the factorization above, we have$\frac{7^{7a} + 1}{7^a + 1} = \big(7^{3a} + 3\cdot7^{2a} + 3\cdot7^a + 1\big)\big(7^{\frac{5a}{2} + \frac{1}{2}} + 7^{\frac{3a}{2} + \frac{1}{2}} + 7^{\frac{a}{2} + \frac{1}{2}}\big)$,

 which is clearly the product of two integers greater than 1, and thus composite.

 Therefore, $\frac{7^{7a} + 1}{7^a + 1}$ is always composite, and hence $7^{7a} + 1$ has at least two more prime factors than $7^a + 1$.

 By the inductive hypothesis, $7^a + 1$ has at least $2n + 3$ prime factors, so $7^{7a} + 1$ has at least $(2n + 3) + 2 = 2n + 5$ prime factors.

 By induction, $7^{7^n} + 1$ is the product of at least $2n + 3$ (not necessarily distinct) primes for all nonnegative integers $n$.  

 
<i>Alternate solutions are always welcome. If you have a different, elegant solution to this problem, please \href{https://artofproblemsolving.com/wiki/index.php?title=2007\_USAMO\_Problems/Problem\_5&action=edit}{add it} to this page.</i>

 


	\item (\textit{USAMO 2007}) Let $ABC$ be an acute triangle with $\omega$, $\Omega$, and $R$ being its incircle, circumcircle, and circumradius, respectively. Circle $\omega_A$ is tangent internally to $\Omega$ at $A$ and tangent externally to $\omega$. Circle $\Omega_A$ is tangent internally to $\Omega$ at $A$ and tangent internally to $\omega$. Let $P_A$ and $Q_A$ denote the centers of $\omega_A$ and $\Omega_A$, respectively.  Define points $P_B$, $Q_B$, $P_C$, $Q_C$ analogously.  Prove that

 \[8P_AQ_A \cdot P_BQ_B \cdot P_CQ_C \le R^3,\]
with equality if and only if triangle $ABC$ is equilateral.


 



\textbf{Solution 1}


                    \begin{figure}[H]
                        \centering
                        \includegraphics[width=0.5\textwidth]{Solution/USAMO_I_2007_6_Solution_Figure_1.png}
                    \end{figure}
                    

 \textbf{Lemma.}
 \[P_{A}Q_{A}=\frac{ 4R^{2}(s-a)^{2}(s-b)(s-c)}{rb^{2}c^{2}}\]

 \textit{Proof.} Note $P_{A}$ and $Q_{A}$ lie on $AO$ since for a pair of tangent circles, the point of tangency and the two centers are \href{https://artofproblemsolving.com/wiki/index.php?title=Collinear}{collinear}. 

 Let $\omega$ touch $BC$, $CA$, and $AB$ at $D$, $E$, and $F$, respectively. Note $AE=AF=s-a$. Consider an \href{https://artofproblemsolving.com/wiki/index.php?title=Inversion}{inversion}, $\mathcal{I}$, centered at $A$, passing through $E$, $F$. Since $IE\perp AE$, $\omega$ is \href{https://artofproblemsolving.com/wiki/index.php?title=Orthogonal&action=edit&redlink=1}{orthogonal} to the inversion circle, so $\mathcal{I}(\omega)=\omega$. Consider $\mathcal{I}(\omega_{A})=\omega_{A}'$. Note that $\omega_{A}$ passes through $A$ and is tangent to $\omega_{A}$, hence $\omega_{A}'$ is a line that is tangent to $\omega$. Furthermore, $\omega_{A}'\perp AO$ because $\omega_{A}$ is symmetric about $OA$, so the inversion preserves that reflective symmetry. Since it is a line that is symmetric about $AO$, it must be perpendicular to $AO$. Likewise, $\mathcal{I}(\Omega_{A})=\Omega_{A}'$ is the other line tangent to $\omega$ and \href{https://artofproblemsolving.com/wiki/index.php?title=Perpendicular}{perpendicular} to $AO$. 

 Let $\omega_{A} \cap AO=X$ and $\omega_{A}' \cap AO=X'$ (second intersection).

 Let $\Omega_{A} \cap AO=Y$ and $\Omega_{A}' \cap AO=Y'$ (second intersection).

 Evidently, $AX=2AP_{A}$ and $AY=2AQ_{A}$. We want:
 \[\star=AQ_{A}-AP_{A}=\frac{1}{2}(AY-AX)=\frac{(s-a)^{2}}{2}\left(\frac{1}{AY'}-\frac{1}{AX'}\right)\]by inversion. Note that $\omega_{A}' || \Omega_{A}'$, and they are tangent to $\omega$, so the \href{https://artofproblemsolving.com/wiki/index.php?title=Distance}{distance} between those lines is $2r=AX'-AY'$. Drop a perpendicular from $I$ to $AO$, touching at $H$. Then $AH=AI\cos\angle OAI=AI\cos\frac{1}{2}|\angle B-\angle C|$. Then $AX', AY' = AI\cos\frac{1}{2}|\angle B-\angle C|\pm r$. So $AX'\cdot AY' = AI^{2}\cos^{2}\frac{1}{2}|\angle B-\angle C|-r^{2}$.
 \begin{align*} \star &= \frac{(s-a)^{2}}{2}\cdot\frac{AX' - AY'}{AY'\cdot AX'} \\ &= \frac{r(s-a)^{2}}{AI^{2}\cos^{2}\frac{1}{2}|\angle B-\angle C|-r^{2}} \\ &= \frac{\frac{(s-a)^{2}}{r}}{\left(\frac{AI}{r}\right)^{2}\cos^{2}\frac{1}{2}|\angle B-\angle C|-1}. \end{align*}
Note that $\frac{AI}{r} = \frac{1}{\sin\frac{A}{2}}$. Applying the double angle formulas and $1-\cos\angle A=\frac{2(s-b)(s-c)}{bc}$, we get
 \begin{align*} \star &= \frac{\frac{(s-a)^{2}}{r}}{\frac{1+\cos (\angle B-\angle C)}{1-\cos A}+1} \\ &= \frac{\frac{(s-a)^{2}}{r}\cdot (1-\cos \angle A)}{\cos(\angle B-\angle C)+\cos(\pi-\angle B-\angle C)} \\ &= \frac{(s-a)^{2}(1-\cos \angle A)}{2r\sin \angle B\sin \angle C} \\ &= \frac{(s-a)^{2}(s-b)(s-c)}{rbc\sin\angle B\sin\angle C} \\ P_{A}Q_{A} &= \frac{4R^{2}(s-a)^{2}(s-b)(s-c)}{rb^{2}c^{2}} \end{align*}


 \textbf{End Lemma}

 The problem becomes:
 \begin{align*} 8\prod \frac{ 4R^{2}(s-a)^{2}(s-b)(s-c)}{rb^{2}c^{2}} &\leq R^{3} \\ \frac{2^{9}R^{6}\left[\prod(s-a)\right]^{4}}{r^{3}a^{4}b^{4}c^{4}} &\leq R^{3} \\ \left(\frac{4AR}{abc}\right)^{4}\cdot\left(\frac{A}{rs}\right)^{4}2rR^{2} &\leq R^{3} \\ 2r &\leq R, \end{align*}
which is true because $OI^{2}=R(R-2r)$, equality is when the circumcenter and incenter coincide. As before, $\angle OAI=\frac{1}{2}|\angle B-\angle C|=0$, so, by symmetry, $\angle A=\angle B=\angle C$. Hence the inequality is true iff $\triangle ABC$ is equilateral.

 \textbf{Comment:} It is much easier to determine $AP_{A}$ by considering $\triangle IAP_{A}$. We have $AI$, $\angle IAO$, $IP_{A}=r+r(P_{A})$, and $AP_{A}=r(P_{A})$. However, the inversion is always nice to use. This also gives an easy construction for $w_{A}$ because the tangency point is collinear with the intersection of $w_{A}'$ and $w$.

 <i>Alternate solutions are always welcome. If you have a different, elegant solution to this problem, please \href{https://artofproblemsolving.com/wiki/index.php?title=2007\_USAMO\_Problems/Problem\_6&action=edit}{add it} to this page.</i>

 


\end{enumerate}